\section{Initial normal recoding of an arbitrary coherent ideal}

\begin{definition}[The whole functional stabilized ideal]
\label{hir:disjoint-coherent-normal-record}
\uses{hir:disjoint-zero-coordinate-initial-record}
\lean{Hironaka.HironakaDisjointCoherentNormalRecoding.baseMarked,Hironaka.HironakaDisjointCoherentNormalRecoding.stabilizedIdeal,Hironaka.HironakaDisjointCoherentNormalRecoding.stabilizedMarked}
\leanok
Let $I$ be a coherent ideal of holomorphic germs on $\mathbf C^n$ and
$M_I=(\mathbf C^n,\mathbf C^n,\varnothing,I,1)$. Write $p(x,t)=x$,
$s(x)=(x,0)$, and $H=\{t=0\}$. Define the whole ideal sheaf
\[
 J(I)=p^*I+\mathcal I_H,
 \qquad M_{J(I)}=(\mathbf C^{n+1},\mathbf C^{n+1},\varnothing,J(I),1).
\]
The sum is the ideal generated by both whole functional stalks at each
actual point. In particular $I$ is retained with its original multiplicities;
no radical or ideal of its support is substituted. Use the established
constructor embedding $\iota_q(v)=(1,\operatorname{inr}((0,v)))$ of
$\operatorname{ValTy}_q$ into $\operatorname{ValTy}_{q+1}$.
\end{definition}

\begin{lemma}[Entire functional section trace]
\label{hir:disjoint-coherent-normal-whole-trace}
\uses{hir:disjoint-coherent-normal-record}
\lean{Hironaka.HironakaDisjointCoherentNormalRecoding.stabilizedIdeal_isCoherent,Hironaka.HironakaDisjointCoherentNormalRecoding.section_pullback_stabilizedIdeal}
\leanok
The stabilized functional ideal is coherent and $s^*J(I)=I$ as whole ideal
sheaves at every original point.
\end{lemma}
\begin{proof}
The plane $H$ is smooth and closed, so $\mathcal I_H$ is coherent. Holomorphic
pullback and finite sum preserve coherence, proving the first assertion.
Pullback distributes over the sum. Since $p\circ s$ is the identity,
$s^*p^*I=I$. Each member germ of $\mathcal I_H$ restricts along $s$ to a germ
identically zero near the base point. Every member of its pulled-back ideal
is a finite holomorphic combination of these restricted germs, so it too
is the zero germ and belongs to $I$. Hence $s^*J(I)\subseteq I$. Conversely,
every member of $I=s^*p^*I$ belongs to the pulled-back sum, proving the reverse
inclusion. This is equality of the original functional ideal predicates,
with no restriction to an origin stalk or to a chosen finite presentation.
\end{proof}

\begin{lemma}[The actual stabilized initial regime]
\label{hir:disjoint-coherent-normal-initial-regime}
\uses{hir:disjoint-coherent-normal-whole-trace,hir:actual-zero-coordinate-genuine-contact-seed}
\lean{Hironaka.HironakaDisjointCoherentNormalRecoding.stabilized_isInvDomain,Hironaka.HironakaDisjointCoherentNormalRecoding.stabilized_cosupport,Hironaka.HironakaDisjointCoherentNormalRecoding.stabilized_order_one,Hironaka.HironakaDisjointCoherentNormalRecoding.stabilized_not_special,Hironaka.HironakaDisjointCoherentNormalRecoding.stabilized_initial_orderMask}
\leanok
Suppose $M_I$ is in the invariant domain of dimension at most $q$.
Then $M_{J(I)}$ is in the domain of dimension at most $q+1$. Its cosupport
is $s(\operatorname{cosupp}M_I)$, its nonmonomial order is exactly one on
that cosupport, its initial order-one mask is the whole ambient manifold,
and its initial level is zero with no special points.
\end{lemma}
\begin{proof}
The coordinate $t$ belongs to the normal summand and has derivative one,
so it is a genuine global maximal-contact direction of the marking-one
stabilized ideal and bounds its order everywhere by one. A cosupport point
must lie in $H$. The entire section-trace equality then identifies its
remaining order-one conditions with those of the original ideal. This
proves the cosupport image statement. The empty boundary identifies
nonmonomial order with ordinary order, so its value on that cosupport is
one and its order-one mask is the whole ambient manifold.
The original domain gives $n\le q$. The stabilized carrier has dimension
$n+1\le q+1$. A base germ nonzero arbitrarily near a projected point pulls
back to a member of $J(I)$ nonzero arbitrarily near the ambient point:
use a parallel section through that point to test every ambient
neighborhood. Coherence and the empty boundary now give the stabilized
domain. There are no original boundary members, so the initial level is
zero. If an initial codimension-one special point existed, maximal contact
by $t$ would force the stabilized cosupport locally to be all of $H$.
Choose an actual original germ nonzero arbitrarily nearby; its pullback
belongs to $J(I)$. Coherence makes this membership valid at every nearby
point. Vanishing on the locally full zero plane and restriction by $s$
would make that original germ identically zero nearby, a contradiction.
Thus the special part is empty as claimed.
\end{proof}

\begin{lemma}[The actual chosen initial contact value]
\label{hir:disjoint-coherent-normal-contact-value}
\uses{hir:disjoint-coherent-normal-initial-regime,hir:disjoint-coherent-normal-whole-trace,hir:actual-closed-section-all-index-initial-values,mc:contact-lower-value}
\lean{Hironaka.HironakaDisjointCoherentNormalRecoding.stabilized_stepOne_initial_value}
\leanok
For $n>0$ and every original cosupport point $x$, the actual Step One
invariant of $M_{J(I)}$, using the depth-$q$ canonical family in ambient
dimension $n+1$, is initially $(0,\operatorname{inv}_q(M_I,x))$.
\end{lemma}
\begin{proof}
For any coherent marking-one ideal, homogenization and each coefficient
operation retain the entire original ideal and mark: the only derivative
term has order zero, exponent one, and mark one. Thus the literal contact
record for $t$ is $(\mathbf C^{n+1},H,\varnothing,J(I),1)$. It lies in the
depth-$q$ domain: $H$ has dimension $n\le q$, and original nonzero germs
pull back to germs nonzero arbitrarily nearby on $H$. Restrict both
ambient manifolds to their whole open sets and apply the existing
initial closed-section value theorem to this actual coherent contact
record. Its entire section trace is $I$, so its value is exactly the
original depth-$q$ value. The nonspecial initial regime internally
supplies the actual chosen lower domain and cosupport. Reference
independence replaces the independently chosen contact direction by $t$.
All towers and tails are empty, so the canonical-tail requirements are
automatic and both run clocks are zero. The whole-open restriction
property identifies the resulting contact value with the preceding
closed-section value. The outer Step One level is zero.
\end{proof}

\begin{theorem}[Nonradical full ambient initial recoding]
\label{hir:disjoint-coherent-normal-initial-value}
\uses{hir:disjoint-coherent-normal-contact-value,hir:disjoint-coherent-normal-record}
\lean{Hironaka.HironakaDisjointCoherentNormalRecoding.stabilized_initial_value}
\leanok
For $n>0$, a natural original domain record $M_I$ of dimension at most $q$,
and $x\in\operatorname{cosupp}M_I$, the actual canonical family satisfies
\[
 \operatorname{inv}_{q+1}(M_{J(I)},s(x))
 =\iota_q\bigl(\operatorname{inv}_q(M_I,x)\bigr).
\]
\end{theorem}
\begin{proof}
The initial general invariant has first component one by the stabilized
order calculation. With empty boundary and mark one, its actual initial
companion on the order mask is restricted-equivalent to the stabilized
record restricted to that mask: the monomial factor is the unit ideal
and the nonmonomial factor is the original ideal restricted to the
carrier. The companion's domain and maximal order follow from the actual
general-step construction. The actual maximal-order Step One invariant
is invariant under this representative change and under restriction to
the order mask. Its finite value is therefore the pair calculated in the
preceding lemma. The definition of the general invariant now gives the
stated finite right-branch expression.
\end{proof}

\begin{theorem}[First maximum locus with all multiplicities retained]
\label{hir:disjoint-coherent-normal-initial-centre}
\uses{hir:disjoint-coherent-normal-initial-value}
\lean{Hironaka.HironakaDisjointCoherentNormalRecoding.stabilized_initial_maximum_locus,Hironaka.HironakaDisjointCoherentNormalRecoding.stabilized_canonical_initial_maximum_locus}
\leanok
The initial depth-$(q+1)$ maximum locus of $M_{J(I)}$ is the zero-section
image of the initial depth-$q$ maximum locus of $M_I$. At $q=n>0$, this
is the exact first-centre formula for the actual canonical invariants in
ambient dimensions $n+1$ and $n$.
\end{theorem}
\begin{proof}
The fixed constructor embedding $\iota_q$ preserves and reflects order.
Every stabilized cosupport point is the section of its projected original
cosupport point. Substitute the proved initial recoding into its defining
comparison with every other cosupport point. Order reflection gives the
original maximum condition, and order preservation gives the converse.
For $q=n>0$ the definition of the canonical invariant is the corresponding
inductive family in both ambient dimensions; substitute these actual
definitions into the maximum-locus identity.
\end{proof}

\begin{lemma}[Empty cosupport is retained]
\label{hir:disjoint-coherent-normal-initial-empty}
\uses{hir:disjoint-coherent-normal-initial-regime}
\lean{Hironaka.HironakaDisjointCoherentNormalRecoding.stabilized_cosupport_eq_empty_iff}
\leanok
The stabilized cosupport is empty exactly when the original cosupport is
empty.
\end{lemma}
\begin{proof}
Each original cosupport point has its section point in the stabilized
cosupport; every stabilized cosupport point projects into the original
one. Hence one cosupport is empty if and only if the other is empty.
The maximum-locus identity consequently includes empty cosupports without
an additional nonemptiness assumption.
\end{proof}

\begin{remark}
Controlled transforms can be nonradical even when the initial ideal comes
from a reduced analytic subset. The whole functional ideal in this
statement retains those multiplicities. The theorem remains an initial,
empty-boundary law. It neither resets an exceptional boundary after a
blowup nor supplies the full historical reference comparison required to
iterate canonical resolution intrinsically.
\end{remark}
